%! TEX root = ../note-dqm-transform.tex

\section{Preservation of DQM by measurable maps}
\label{sec--dqm-and-meas-maps}

Suppose now that instead of observing \(X \sim Q_{0}\), we observe the
transformation \(T (X)\), where \(T : (\mathbf{X}, \mathcal{B}) \to \left(
\mathbf{X}^{\prime}, \mathcal{B}^{\prime} \right)\) is a measurable
transformation.
The associated model is then the collection of pushforward probability measures:
\begin{equation}
  \mathbf{P} = \{ P_{\theta} : \theta \in \Theta\} \quad \text{where} \quad
  P_{\theta} = Q_{\theta} \circ T^{- 1}.
  \label{eqn--P-T-def}
\end{equation}
As before,
\begin{equation}
  P_{\theta} =  P_{\theta, 0}^{\perp} + P_{\theta, 0}
  \quad \text{where} \quad
  P_{\theta, 0}^{\perp} \perp P_{0},
  \quad \text{and} \quad
  P_{\theta, 0} \ll P_{0}.
  \label{eqn--Ptheta-leb-decomp}
\end{equation}
The question is then whether \eqref{eqn--dqm-def} holds for the quantities in
\eqref{eqn--Ptheta-leb-decomp}.
An intuitive first step would be to inspect if a relationship can be established
between the decompositon in \eqref{eqn--Ptheta-leb-decomp}
and a ``Lebesgue pseudo-decompositon'':
\begin{equation}
  P_{\theta} = \tilde{P}_{\theta, 0}^{\perp} + \tilde{P}_{\theta, 0}
  \quad \text{where} \quad
  \tilde{P}_{\theta, 0}^{\perp} = Q_{\theta, 0}^{\perp} \circ T^{- 1},
  \quad \text{and} \quad
  \tilde{P}_{\theta, 0} = Q_{\theta, 0} \circ T^{- 1}.
  \label{eqn--Ptheta-leb-decomp-pseudo}
\end{equation}
The decomposition \eqref{eqn--Ptheta-leb-decomp-pseudo} is called a ``Lebesgue
pseudo-decomposition'' since it is not true in general that \(\tilde{P}_{\theta,
0}^{\perp} = P_{\theta, 0}^{\perp}\) and \(\tilde{P}_{\theta, 0} = P_{\theta,
0}\).
In fact, only the inequalities
\(\tilde{P}_{\theta, 0} \leq P_{\theta, 0}\) and
\(\tilde{P}_{\theta, 0}^{\perp} \geq P_{\theta, 0}^{\perp}\) hold.

\begin{lemma}
\label{lem--density-preservation}
Let \(Q\) be a measure and \(\mu\) be a probability measure on \((\mathbf{X},
\mathcal{B})\) with \(Q \ll \mu\).
Let \(T : \mathbf{X} \mapsto \mathbf{X}^{\prime}\) be measurable from
\((\mathbf{X}, \mathcal{B})\) to \(\left( \mathbf{X}^{\prime},
\mathcal{B}^{\prime} \right)\), and denote \(\nu := \mu \circ T^{- 1}\).
Then \(P = Q \circ T^{- 1} \ll \nu\) and
\begin{equation}
  \frac{\mathrm{d} P}{\mathrm{d} \nu} (t) = \mu \left[ \frac{\mathrm{d}
  Q}{\mathrm{d} \mu} \middle| T = t \right] \quad \nu \text{ almost everywhere}.
  \label{eqn--density-preservation}
\end{equation}
\end{lemma}

\begin{proof}
See \Cref{sec--prf--lem--density-preservation}.
\end{proof}

\begin{lemma}
\label{lem--leb-decomp-pseudo-ineqs}
Let \(Q\) be a measure and \(\mu\) be positive measures on \((\mathbf{X},
\mathcal{B})\).
Denote the Lebesgue decomposition of \(Q\) in terms of \(\mu\) by
\begin{equation*}
  Q = Q^{\perp} + Q^{\mathrm{ac}} \quad \text{where} \quad Q^{\perp} \perp \mu
  \quad \text{and} \quad Q^{\mathrm{ac}} \ll \mu.
\end{equation*}
Let \(T : \mathbf{X} \mapsto \mathbf{X}^{\prime}\) be measurable from
\((\mathbf{X}, \mathcal{B})\) to \(\left( \mathbf{X}^{\prime},
\mathcal{B}^{\prime} \right)\), and denote \(\nu := \mu \circ T^{- 1}\) and
\(P = Q \circ T^{- 1}\).
Denote the following:
\begin{equation*}
  \begin{gathered}
    P = P^{\perp} + P^{\mathrm{ac}} \quad \text{where} \quad P^{\perp} \perp \nu
    \quad \text{and} \quad P^{\mathrm{ac}} \ll \nu, \\
    P = \tilde{P}^{\perp} + \tilde{P}^{\mathrm{ac}}
    \quad \text{where} \quad
    \tilde{P}^{\perp} = Q^{\perp} \circ T^{- 1} \quad \text{and} \quad
    \tilde{P}^{\mathrm{ac}} = Q^{\mathrm{ac}} \circ T^{- 1}.
  \end{gathered}
\end{equation*}
Then \(\tilde{P}^{\mathrm{ac}} \leq P^{\mathrm{ac}}\) and \(\tilde{P}^{\perp}
\geq P^{\perp}\) hold.
\end{lemma}

\begin{proof}[Proof of \Cref{lem--leb-decomp-pseudo-ineqs}]
By \Cref{lem--density-preservation}, \(\tilde{P}^{\mathrm{ac}} \ll \nu\).
Therefore, \(\tilde{P}^{\mathrm{ac}}\) has no singular component with respect to
\(\nu\).
Now take the Lebesgue decomposition of \(\tilde{P}^{\perp}\) with respect to
\(\nu\):
\begin{equation*}
  \tilde{P}^{\perp} = \hat{P}^{\perp} + \hat{P}^{\mathrm{ac}} \quad \text{where}
  \quad \hat{P}^{\perp} \perp \nu \quad \text{and} \quad \hat{P}^{\mathrm{ac}}
  \ll \nu.
\end{equation*}
But then, \(P = \hat{P}^{\perp} + \hat{P}^{\mathrm{ac}} +
\tilde{P}^{\mathrm{ac}}, \quad \hat{P}^{\perp} \perp \nu \quad \text{and} \quad
\hat{P}^{\mathrm{ac}} + \tilde{P}^{\mathrm{ac}} \ll \nu\).
By uniqueness of the Lebesgue decomposition, it is necessary that
\(\hat{P}^{\perp} = P^{\perp}\).
Thus \(\tilde{P}^{\perp} = P^{\perp} + \hat{P}^{\mathrm{ac}} \geq P^{\perp}\).
\end{proof}

\begin{theorem}
\label{thm--P-dqm}
Let \((\mathbf{X}, \mathcal{B}, \mu)\) be a probability space, \(\left(
\mathbf{X}^{\prime}, \mathcal{B}^{\prime} \right)\) be a measurable space,
\(T : \mathbf{X} \to \mathbf{X}^{\prime}\) be a measurable function.
Let \(\mathbf{Q}\) be a probability model and \(\mathbf{P}\) its pusforward
family under \(T\).
If \(\mathbf{Q}\) is DQM at \(0 \in \Theta\) with score
\(\dot{\ell}_{0}\), then \(\mathbf{P}\) is DQM at \(0\) with score
\begin{equation}
  \dot{\ell}_{T, 0} (t) = Q_{0} \left[ \dot{\ell}_{0} \middle| T = t \right].
  \label{eqn--dotl-marg-T-def}
\end{equation}
\end{theorem}

\begin{proof}[Proof of \Cref{thm--P-dqm}]
By \Cref{lem--leb-decomp-pseudo-ineqs},
\(P_{\theta, 0}^{\perp} \left( \mathbf{X}^{\prime} \right) \leq
\tilde{P}_{\theta, 0}^{\perp} \left( \mathbf{X}^{\prime} \right) =
Q_{\theta, 0}^{\perp} (\mathbf{X}) = o \left( |\theta|^{2} \right)\), by
\Cref{def--dqm} \ref{def--dqm-sing} applied to \(\mathbf{Q}\) at 0.
Thus, \(P_{\theta, 0}^{\perp}\) satisfies
\Cref{def--dqm} \ref{def--dqm-sing}.

Let \(p_{\theta, 0} = \mathrm{d} P_{\theta, 0} / \mathrm{d} P_{0}\) and
\(\tilde{p}_{\theta, 0} = \mathrm{d} \tilde{P}_{\theta, 0} / \mathrm{d} P_{0}\).
By \Cref{lem--leb-decomp-pseudo-ineqs}, since \(\tilde{P}_{\theta, 0} \leq
P_{\theta, 0}\) everywhere, \(\tilde{p}_{\theta, 0} \leq p_{\theta, 0}\) almost
everywhere \(P_{0}\).
Let
\begin{equation*}
  \psi_{\theta, 0} = p_{\theta, 0} - \tilde{p}_{\theta, 0},
  \quad \text{and} \quad
  \varphi_{\theta, 0} = \sqrt{p_{\theta, 0}} - \sqrt{\tilde{p}_{\theta, 0}}.
\end{equation*}
Note that
\begin{equation*}
    \left| \varphi_{\theta, 0} \right|^{2} = \left| \sqrt{p_{\theta, 0}} -
    \sqrt{\tilde{p}_{\theta, 0}} \right|^{2} \leq \left| \sqrt{p_{\theta, 0}} -
    \sqrt{\tilde{p}_{\theta, 0}} \right| \left| \sqrt{p_{\theta, 0}} +
    \sqrt{\tilde{p}_{\theta, 0}} \right| = \left| p_{\theta, 0} -
    \tilde{p}_{\theta, 0} \right| = \psi_{\theta, 0}.
\end{equation*}
Then
\begin{equation}
  \begin{split}
    P_{0} \left[ \varphi_{\theta, 0}^{2} \right] \leq
    & \,
    P_{0} \psi_{\theta, 0} =
    P_{\theta, 0} \left( \mathbf{X}^{\prime} \right) -
    \tilde{P}_{\theta, 0} \left( \mathbf{X}^{\prime} \right) \leq P_{\theta}
    \left( \mathbf{X}^{\prime} \right) - \tilde{P}_{\theta, 0} \left(
    \mathbf{X}^{\prime} \right) = \tilde{P}_{\theta, 0}^{\perp} \left(
    \mathbf{X}^{\prime} \right) \\
    =
    & \, Q_{\theta, 0}^{\perp} (\mathbf{X})
    = o \left( |\theta|^{2} \right).
  \end{split}
  \label{eqn--ac-part-rem-neg}
\end{equation}
Define
\begin{equation}
  \begin{gathered}
    \rho_{\theta, 0} = q_{\theta, 0} - 1 - \frac{1}{2} \dot{\ell}_{0}^{\prime}
    \theta, \qquad r_{\theta, 0} = p_{\theta, 0} - 1 - \frac{1}{2}
    \dot{\ell}_{T, 0}^{\prime} \theta \quad \text{and} \quad \tilde{r}_{\theta,
    0} = \tilde{p}_{\theta, 0} - 1 - \frac{1}{2} \dot{\ell}_{T, 0}^{\prime}
    \theta, \\
    \text{so that} \quad r_{\theta, 0} = \tilde{r}_{\theta, 0} +
    \varphi_{\theta, 0}.
  \end{gathered}
  \label{eqn--dqm-P-rem-defs}
\end{equation}
From \eqref{eqn--ac-part-rem-neg}, \eqref{eqn--dqm-P-rem-defs} and the fact that
\(P_{0} \left[ \rho_{\theta, 0}^{2} \right] \leq 2 P_{0} \left[
\tilde{\rho}_{\theta, 0}^{2} \right] + 2 P_{0} \left[ \varphi_{\theta, 0}^{2}
\right]\), it remains to show then that \(P_{0} \left[
\tilde{\rho}_{\theta, 0}^{2} \right] = o \left( |\theta|^{2} \right)\).

Since \(\widetilde{P}_{\theta, 0} = Q_{\theta, 0} \circ T^{- 1}\), \(P_{0} =
Q_{0} \circ T^{- 1}\) and \(Q_{\theta, 0} \ll Q_{0}\), it follows from
\Cref{lem--density-preservation} that \(\tilde{p}_{\theta, 0} (t) = Q_{0} \left[
q_{\theta, 0} \middle| T = t \right]\).
Since \(\dot{\ell}_{T, 0} (t) = Q_{0} \left[ \dot{\ell}_{0} \middle| T = t
\right]\), \(\tilde{r}_{\theta} = Q_{0} \left[ \rho_{\theta, 0} \middle| T = t
\right]\).
By Jensen's inequality, \(\tilde{r}_{\theta, 0}^{2} \leq Q_{0} \left[
\rho_{\theta, 0}^{2} \middle| T = t \right]\).
By \eqref{eqn--dqm-def} in \Cref{def--dqm} \ref{def--dqm-ac-diff} and the Law of
Iterated Expectations,
\(Q_{0} \left[ \tilde{r}_{\theta, 0}^{2} \right] \leq Q_{0} \left[
\rho_{\theta, 0}^{2} \right] = o \left( |\theta|^{2} \right)\).
Therefore, \(\mathbf{P}\) satisfies \Cref{def--dqm} \ref{def--dqm-ac-diff} at 0
with the prescribed score function.
\end{proof}

%%% Local Variables:
%%% mode: LaTeX
%%% TeX-master: "../note-dqm-transform"
%%% End:
